\chapter{Applications conformes et théorème de l'application de Riemann}
\label{ch:b3:conformal}

Une bijection \omterm{def:b3:holomorphic:holo}{holomorphe} entre deux domaines transporte toute
l'analyse complexe de l'un à l'autre~: de telles applications ---
\emph{conformes}, car elles préservent les angles --- sont les
isomorphismes du monde \omterm{def:b3:holomorphic:holo}{holomorphe}. Ce chapitre les classe là où
la classification est possible (le disque, le plan~: le lemme de
Schwarz est la clé, une inégalité d'une puissance étonnante),
construit la théorie de compacité des familles \omterm{def:b3:holomorphic:holo}{holomorphes}
(Montel), et prouve le plus profond théorème d'existence du
sujet~: \emph{tout} sous-domaine propre \omterm{def:b3:conformal:simplyconnected}{simplement connexe} de
$\C$, aussi déchiqueté soit son bord, est conformément
équivalent au disque unité. Nous clôturons avec les \omterm{prop:b3:conformal:harmonicholo}{fonctions
harmoniques} et le \omterm{thm:b3:conformal:poisson}{noyau de Poisson}, en résolvant le problème de
Dirichlet sur le disque --- la récompense analytique de la
géométrie conforme. Tout au long, $\mathbb D = D(0,1)$ et
$\mathbb H = \{\operatorname{Im}z > 0\}$.

\section{Applications conformes~; transformations de Möbius}

\begin{definition}\label{def:b3:conformal:conformal}
Une \emph{application conforme}\index{application conforme} (ou
biholomorphisme) entre \omterm{def:b3:topology:topology}{ouverts} est une bijection \omterm{def:b3:holomorphic:holo}{holomorphe}~;
son inverse est automatiquement \omterm{def:b3:holomorphic:holo}{holomorphe}
(le~\cref{cor:b3:residues:openmapping}). Deux domaines sont
\emph{conformément équivalents} si une telle application
existe~; $\operatorname{Aut}(\Omega)$ désigne le groupe des
auto-applications conformes. Là où $f' \neq 0$ --- partout,
pour $f$ injective (preuve
du~\cref{cor:b3:residues:openmapping}) --- la différentielle
est la multiplication par $f'(z) \neq 0$~: une similitude,
donc les applications conformes préservent les angles entre
courbes, y compris l'orientation.
\end{definition}

\begin{example}[Transformations de Möbius]
\label{ex:b3:conformal:mobius}
Pour $\bigl(\begin{smallmatrix}a & b\\ c &
d\end{smallmatrix}\bigr) \in GL_2(\C)$, la \emph{transformation
de Möbius}\index{transformation de Möbius} $z \mapsto
\frac{az + b}{cz + d}$ est conforme de $\C\setminus\{-d/c\}$
sur $\C\setminus\{a/c\}$ (inverse du même type, à partir de
la matrice inverse~; la composition correspond au produit
matriciel). L'\emph{application de Cayley}\index{transformation
de Cayley}
\[
\varphi(z) = \frac{z - \iu}{z + \iu}
\]
envoie $\mathbb H$ conformément sur $\mathbb D$~: en effet
$\abs{z - \iu} < \abs{z + \iu}$ exactement quand $z$ est plus
proche de $\iu$ que de $-\iu$, c'est-à-dire
$\operatorname{Im}z > 0$~; l'inverse est $w \mapsto
\iu\frac{1 + w}{1 - w}$. Les applications de Möbius envoient
la famille cercles-et-droites sur elle-même
(l'\cref{exo:b3:conformal:1}).
\end{example}

\begin{example}[L'application de Joukowski]\label{ex:b3:conformal:joukowski}
Au-delà de Möbius, l'\omterm{def:b3:conformal:conformal}{application conforme} la plus utile des
mathématiques appliquées classiques est
\[
J(z) = \frac12\Bigl(z + \frac1z\Bigr) .
\]
Sur l'extérieur $\Omega = \{\abs z > 1\}$ du disque unité,
$J$ est injective~: $J(z) = J(w)$ donne $(z - w)(1 -
\frac1{zw}) = 0$ et $\abs{zw} > 1$. Sa dérivée $J'(z) =
\frac12(1 - z^{-2})$ s'annule seulement en $z = \pm1$, sur le
bord~: $J$ est une équivalence conforme de $\Omega$ sur son
image, qui est $\C\setminus\intcc{-1}1$ --- le cercle unité
lui-même est plié deux-à-un sur le segment
($J(\eu^{\iu\theta}) = \cos\theta$). Ainsi l'extérieur d'un
\emph{segment}, un plan fendu sans bord lisse, est
conformément l'extérieur d'un disque~: les coins ne sont pas
un obstacle à l'équivalence conforme, seulement à la
régularité du bord. Les images de cercles passant par $\pm1$
mais décentrés sont des courbes en forme de profil d'aile, et
composer $J$ avec des applications de Möbius transportait
l'écoulement autour d'un cylindre --- calculable à la main ---
vers l'écoulement autour d'une aile~: pour la première moitié
du vingtième siècle, cet exemple \emph{était}
l'aérodynamique. C'est aussi la porte vers Tchebychev~: $J$
conjugue $z \mapsto z^n$ au polynôme de Tchebychev $T_n$
(le~\cref{pb:b3:hilbert:1}, partie~V), puisque $J(z^n) =
\cos(n\theta)$ quand $z = \eu^{\iu\theta}$.
\end{example}

\section{Lemme de Schwarz et groupes d'automorphismes}

\begin{theorem}[Lemme de Schwarz]\label{thm:b3:conformal:schwarz}
Soit $f \colon \mathbb D \to \mathbb D$ \omterm{def:b3:holomorphic:holo}{holomorphe} avec
$f(0) = 0$. Alors\index{lemme de Schwarz}
\[
\abs{f(z)} \leq \abs z \ \ (z \in \mathbb D)
\qquad\text{et}\qquad \abs{f'(0)} \leq 1 ,
\]
et si $\abs{f(z_0)} = \abs{z_0}$ pour un $z_0 \neq 0$, ou
$\abs{f'(0)} = 1$, alors $f(z) = \eu^{\iu\theta}z$ est une
rotation.
\end{theorem}

\begin{proof}
$g(z) = f(z)/z$ s'étend holomorphiquement à $\mathbb D$ (la
singularité en $0$ est éliminable~: $g$ est bornée près de
$0$, le~\cref{thm:b3:residues:riemanncw}~; sa valeur en $0$
est $f'(0)$). Sur $\abs z = r < 1$~: $\abs g \leq \frac1r$,
donc par le principe du maximum
(le~\cref{thm:b3:holomorphic:maximum}) $\abs g \leq \frac1r$
sur $\bar D(0,r)$~; faire $r \to 1$~: $\abs g \leq 1$ sur
$\mathbb D$, ce qui est les deux inégalités. L'égalité en un
point intérieur rend $\abs g$ atteignant un maximum
intérieur~: $g$ constante de \omterm{def:b3:modules:module}{module} $1$.
\end{proof}

\begin{theorem}[Automorphismes du disque]
\label{thm:b3:conformal:autdisc}
Pour $a \in \mathbb D$, le \emph{facteur de
Blaschke}\index{facteur de Blaschke}
\[
\varphi_a(z) = \frac{z - a}{1 - \bar a z}
\]
est un automorphisme de $\mathbb D$ échangeant $a$ et $0$,
avec $\varphi_a^{-1} = \varphi_{-a}$. Tout automorphisme de
$\mathbb D$ est $\eu^{\iu\theta}\varphi_a$ pour un unique
$\theta \in \R/2\pi\Z$, $a \in \mathbb D$.
\end{theorem}

\begin{proof}
Sur $\abs z = 1$~: $\abs{1 - \bar az} = \abs{\bar z}\abs{1 -
\bar az} = \abs{\bar z - \bar a\abs z^2} = \abs{\bar z - \bar
a} = \abs{z - a}$, donc $\abs{\varphi_a} = 1$ là~; par le
principe du maximum $\varphi_a(\mathbb D) \subseteq
\bar{\mathbb D}$, et l'ouverture met l'image dans $\mathbb
D$. L'identité \omterm{def:b3:galois:algebraic}{algébrique} $\varphi_{-a}\circ\varphi_a =
\mathrm{id}$ (calcul direct) montre la bijectivité. Soit
maintenant $f \in \operatorname{Aut}(\mathbb D)$ et $a =
f^{-1}(0)$~: $g = f\circ\varphi_{-a}$ est un automorphisme
fixant $0$. Schwarz appliqué à $g$ et à $g^{-1}$~:
$\abs{g(z)} \leq \abs z$ et $\abs{g^{-1}(w)} \leq \abs w$,
donc $\abs{g(z)} = \abs z$~: rotation, $g =
\eu^{\iu\theta}\,\mathrm{id}$, c'est-à-dire $f =
\eu^{\iu\theta}\varphi_a$. Unicité~: $a = f^{-1}(0)$ et
$\theta$ par évaluation de type $f'$ (ou de $f(0) =
-\eu^{\iu\theta}a$ et une valeur de plus).
\end{proof}

\begin{theorem}[Automorphismes du plan]
\label{thm:b3:conformal:autplane}
$\operatorname{Aut}(\C) = \{z \mapsto az + b : a \in \C^*,\
b \in \C\}$. Par conséquent $\C$ et $\mathbb D$ ne sont pas
conformément équivalents.
\end{theorem}

\begin{proof}
Soit $f \in \operatorname{Aut}(\C)$ et considérons $g(z) =
f(1/z)$ sur $\C^*$~: une \omterm{def:b3:holomorphic:holo}{fonction holomorphe} avec une
\omterm{def:b3:residues:singularities}{singularité isolée} en $0$. Si elle était essentielle,
Casorati--Weierstrass (le~\cref{thm:b3:residues:riemanncw})
rendrait $g\bigl(D(0,\varepsilon)\setminus\{0\}\bigr)$ \omterm{def:b3:topology:interior}{dense},
tandis que $f(D(0, 1))$ est \omterm{def:b3:topology:topology}{ouvert} et disjoint de celle-ci
($f$ injective~: les deux ensembles sont images d'ensembles
disjoints) --- impossible pour un ensemble \omterm{def:b3:topology:interior}{dense} et un
\omterm{def:b3:topology:topology}{ouvert} non vide. Donc $0$ est un pôle ou éliminable pour
$g$, c'est-à-dire $\abs{f(z)}$ a au plus une croissance
polynomiale~: $f$ est un polynôme
(l'\cref{exo:b3:holomorphic:4}(a)). L'injectivité force le
\omterm{def:b3:galois:extension}{degré} $1$~: un polynôme de \omterm{def:b3:galois:extension}{degré} plus élevé a soit une
racine multiple de $f - c$ quelque part ($f'$ s'annule) soit
plusieurs préimages distinctes (d'Alembert--Gauss,
le~\cref{pb:b3:holomorphic:1})~; dans les deux cas
l'injectivité échoue. Enfin, une \omterm{def:b3:conformal:conformal}{application conforme} $\C
\to \mathbb D$ serait une fonction entière bornée~:
constante (Liouville) --- pas d'équivalence.
\end{proof}

\section{Théorème de Montel}

\begin{theorem}[Montel]\label{thm:b3:conformal:montel}
Soit $\mathcal F \subseteq \mathcal H(\Omega)$
\emph{localement bornée}~: tout point a un \omterm{def:b3:topology:topology}{voisinage} sur
lequel $\sup_{f\in \mathcal F}\sup\abs f < \infty$. Alors
toute suite de $\mathcal F$ a une sous-suite convergeant
uniformément sur tous les sous-ensembles compacts de
$\Omega$ (vers une limite
\omterm{def:b3:holomorphic:holo}{holomorphe}).\index{théorème de Montel}
\end{theorem}

\begin{proof}
\omterm{def:b3:complete:equicontinuous}{Équicontinuité} locale~: si $\abs f \leq M$ sur $D(a, 2r)
\subseteq \Omega$ pour tout $f \in \mathcal F$, la formule de
Cauchy donne, pour $z, z' \in D(a, r)$,
\[
\abs{f(z) - f(z')} = \frac{\abs{z - z'}}{2\pi}
\Bigl|\int_{C_{2r}}\frac{f(w)\,\dd w}{(w-z)(w-z')}\Bigr|
\leq \frac{\abs{z - z'}\,2\pi\cdot2r\,M}{2\pi\,r^2}
= \frac{2M}{r}\,\abs{z - z'} :
\]
une borne de Lipschitz uniforme. Épuiser $\Omega$ par des
compacts $K_m$~; chaque $K_m$ est recouvert par un nombre
fini de tels disques, donc $\mathcal F$ est uniformément
bornée et \omterm{def:b3:complete:equicontinuous}{équicontinue} sur $K_m$~: Arzelà--Ascoli
(le~\cref{thm:b3:complete:ascoli}) extrait une sous-suite
convergeant uniformément sur $K_m$~; diagonaliser sur $m$.
La limite est \omterm{def:b3:holomorphic:holo}{holomorphe} par
le~\cref{thm:b3:holomorphic:weierstrassconv}.
\end{proof}

\section{Le théorème de l'application de Riemann}

\begin{definition}\label{def:b3:conformal:simplyconnected}
Un \omterm{def:b3:topology:topology}{ouvert} \omterm{def:b3:topology:connected}{connexe} $\Omega \subseteq \C$ est \emph{simplement
connexe}\index{domaine simplement connexe} (au sens
homologique, suffisant pour tous nos besoins) si
$\operatorname{Ind}_\gamma(w) = 0$ pour tout cycle $\gamma$
dans $\Omega$ et tout $w \notin \Omega$ --- «~aucun cycle de
$\Omega$ n'entoure un trou~». Par le théorème de Cauchy
global (le~\cref{thm:b3:residues:globalcauchy}) et
la~\cref{prop:b3:holomorphic:primitive}, sur un tel $\Omega$
\emph{toute \omterm{def:b3:holomorphic:holo}{fonction holomorphe} a une primitive}~; d'où
toute $f \in \mathcal H(\Omega)$ sans zéro a un logarithme
\omterm{def:b3:holomorphic:holo}{holomorphe} ($\exp\circ$(primitive de $f'/f$), ajusté d'une
constante, car $(f\eu^{-L})' = 0$) et des racines $n$-ièmes
\omterm{def:b3:holomorphic:holo}{holomorphes} $\eu^{L/n}$.\index{logarithme holomorphe}
\end{definition}

\begin{theorem}[Théorème de l'application de Riemann]
\label{thm:b3:conformal:rmt}
Tout \omterm{def:b3:topology:topology}{ouvert} \omterm{def:b3:conformal:simplyconnected}{simplement connexe} $\Omega \subsetneq \C$,
$\Omega \neq \varnothing$, est conformément équivalent à
$\mathbb D$~; donné $z_0 \in \Omega$, il existe une unique
$f \colon \Omega \to \mathbb D$ conforme avec $f(z_0) = 0$
et $f'(z_0) > 0$.\index{théorème de l'application de Riemann}
\end{theorem}

\begin{proof}
\emph{Étape~0~: la famille est non vide.} Prendre $b \notin
\Omega$~: $z - b$ est sans zéro sur $\Omega$, donc a une
racine carrée \omterm{def:b3:holomorphic:holo}{holomorphe} $h$ ($h^2 = z - b$). $h$ est
injective ($h(z) = h(z')$ élève au carré en $z = z'$), et si
$w \in h(\Omega)$ alors $-w \notin h(\Omega)$ ($h(z) =
-h(z')$ élève aussi au carré en $z = z'$, donnant $w = -w =
0$, impossible car $h$ est sans zéro). Comme $h(\Omega)$ est
\omterm{def:b3:topology:topology}{ouvert}, il contient un disque $D(h(z_0), \rho)$~; alors
$D(-h(z_0), \rho) \cap h(\Omega) = \varnothing$,
c'est-à-dire $\abs{h(z) + h(z_0)} \geq \rho$ pour tout $z
\in \Omega$. D'où
\[
g(z) = \frac{\rho}{2\,\bigl(h(z) + h(z_0)\bigr)}
\]
est \omterm{def:b3:holomorphic:holo}{holomorphe}, injective (une application de Möbius composée
avec l'injective $h$), avec $\abs g \leq \frac12 < 1$. En
composant avec un \omterm{thm:b3:conformal:autdisc}{facteur de Blaschke}
(le~\cref{thm:b3:conformal:autdisc}) pour amener $g(z_0)$ en
$0$, la famille
\[
\mathcal F = \{f \colon \Omega \to \mathbb D \text{
holomorphe, injective, } f(z_0) = 0\}
\]
est non vide.

\emph{Étape~1~: un élément extrémal.} Soit $s =
\sup_{\mathcal F}\abs{f'(z_0)} \in \intoc0{+\infty}$ ($> 0$~:
les membres sont injectifs, donc $f'(z_0) \neq 0$). Prendre
$f_n \in \mathcal F$ avec $\abs{f_n'(z_0)} \to s$~: la
famille est bornée par $1$, donc Montel
(le~\cref{thm:b3:conformal:montel}) extrait $f_n \to f$
uniformément sur les compacts~; $f$ est \omterm{def:b3:holomorphic:holo}{holomorphe}, $f(z_0)
= 0$, $\abs{f'(z_0)} = s$
(le~\cref{thm:b3:holomorphic:weierstrassconv} pour les
dérivées), en particulier $f$ non constante~; $f$ est
injective par Hurwitz (l'\cref{exo:b3:residues:8}(b)), et
$f(\Omega) \subseteq \bar{\mathbb D}$, d'où $\subseteq
\mathbb D$ (application ouverte). Donc $f \in \mathcal F$
atteint le supremum~: $s < \infty$.

\emph{Étape~2~: l'application extrémale est surjective.}
Supposons $a \in \mathbb D\setminus f(\Omega)$. Le transport
de Blaschke $\varphi_a\circ f$ est sans zéro sur le
\omterm{def:b3:conformal:simplyconnected}{simplement connexe} $\Omega$~: il a une racine carrée
\omterm{def:b3:holomorphic:holo}{holomorphe} $F$ (avec $F(\Omega) \subseteq \mathbb D$, car
$\abs F^2 = \abs{\varphi_a\circ f} < 1$), injective (les
carrés distinguent). Normaliser~: $G =
\varphi_{F(z_0)}\circ F \in \mathcal F$. En défaisant~: $f
= \varphi_{-a}\circ s_2 \circ \varphi_{-F(z_0)}\circ G$ où
$s_2(w) = w^2$~; l'application $\Psi =
\varphi_{-a}\circ s_2\circ\varphi_{-F(z_0)} \colon \mathbb D
\to \mathbb D$ est \omterm{def:b3:holomorphic:holo}{holomorphe} avec $\Psi(0) = f(z_0) = 0$ et
n'est \emph{pas} une rotation (elle n'est pas injective~:
$s_2$ ne l'est pas). Lemme de Schwarz (cas strict)~:
$\abs{\Psi'(0)} < 1$, et la règle de la chaîne $f = \Psi\circ
G$ donne $\abs{f'(z_0)} = \abs{\Psi'(0)}\,\abs{G'(z_0)} <
\abs{G'(z_0)}$ --- contredisant la maximalité (noter $G \in
\mathcal F$). D'où $f$ est surjective~: une équivalence
conforme.

\emph{Étape~3~: normalisation et unicité.} Multiplier $f$ par
$\eu^{-\iu\arg f'(z_0)}$ pour rendre $f'(z_0) > 0$ (cela reste
dans $\mathcal F$). Si $f_1, f_2$ marchent toutes deux, $\psi
= f_2\circ f_1^{-1} \in \operatorname{Aut}(\mathbb D)$ fixe
$0$ avec $\psi'(0) = f_2'(z_0)/f_1'(z_0) > 0$~; par
le~\cref{thm:b3:conformal:autdisc} $\psi$ est une rotation
$\eu^{\iu\theta}$ avec $\eu^{\iu\theta} > 0$~: $\psi =
\mathrm{id}$.
\end{proof}

\begin{remark}\label{rem:b3:conformal:rmtscope}
Le théorème est un pur énoncé d'existence d'une portée
étonnante~: un carré, un demi-plan, le complémentaire d'une
fente, la région entre deux cercles tangents, un domaine à
bord fractal --- tous conformément identiques à $\mathbb D$.
Ce qu'il ne donne pas~: aucune formule (les applications
explicites sont l'exception~: l'\cref{exo:b3:conformal:5}), de
comportement au bord (une théorie plus profonde --- le
théorème de Carathéodory --- le traite), ni d'unicité
d'extension à $\C$ ou aux domaines multiplement \omterm{def:b3:topology:connected}{connexes}~:
l'anneau $\{1 < \abs z < 2\}$ n'est \emph{pas}
conformément un disque percé, et des anneaux de rapports de
rayons différents sont inéquivalents (un fait
authentiquement plus dur).
\end{remark}

\section{Fonctions harmoniques et noyau de Poisson}

\begin{proposition}\label{prop:b3:conformal:harmonicholo}
Soit $\Omega$ \omterm{def:b3:conformal:simplyconnected}{simplement connexe} et $u \colon \Omega \to \R$
harmonique\index{fonction harmonique} ($\mathcal C^2$ avec
$\Delta u = u_{xx} + u_{yy} = 0$). Alors $u =
\operatorname{Re}F$ pour une $F$ \omterm{def:b3:holomorphic:holo}{holomorphe}, unique à une
constante imaginaire près. Par conséquent $u$ est $\mathcal
C^\infty$, satisfait la propriété de la moyenne, et obéit au
principe du maximum (pas d'extrémum intérieur strict sauf si
constante).
\end{proposition}

\begin{proof}
$g = u_x - \iu u_y$ satisfait les \omterm{prop:b3:holomorphic:cauchyriemann}{équations de
Cauchy--Riemann} ($P = u_x$, $Q = -u_y$~: $P_x = u_{xx} =
-u_{yy} = Q_y$ et $P_y = u_{xy} = u_{yx} = -Q_x$) avec
dérivées partielles \omterm{def:b3:topology:continuity}{continues}~: $g$ est \omterm{def:b3:holomorphic:holo}{holomorphe}
(la~\cref{prop:b3:holomorphic:cauchyriemann}~; la
$\R$-différentiabilité suit de $\mathcal C^1$). Soit $F_0$
une primitive (\omterm{def:b3:groups:simple}{simple} connexité,
la~\cref{def:b3:conformal:simplyconnected})~; alors
$\operatorname{Re}F_0$ a pour gradient $(u_x, u_y)$
($F_0' = g$ se déplie en exactement cela via Cauchy--Riemann
pour $F_0$), donc $u - \operatorname{Re}F_0$ est constante
($\Omega$ \omterm{def:b3:topology:connected}{connexe})~: ajuster $F = F_0 + c$. Les propriétés
se transfèrent du~\cref{thm:b3:holomorphic:maximum} et de
l'\cref{exo:b3:holomorphic:10} (pour le principe du maximum
appliqué à $u$ elle-même, utiliser $\eu^{F}$ comme là).
\end{proof}

\begin{theorem}[Formule de Poisson~; problème de Dirichlet sur
le disque]\label{thm:b3:conformal:poisson}
Pour $0 \leq r < 1$ définir le \emph{noyau de
Poisson}\index{noyau de Poisson}
\[
P_r(\theta) = \sum_{n\in\Z}r^{\abs n}\eu^{\iu n\theta}
= \frac{1 - r^2}{1 - 2r\cos\theta + r^2} \;>\; 0 .
\]
Soit $g \colon \partial\mathbb D \to \R$ \omterm{def:b3:topology:continuity}{continue}, et poser,
pour $z = r\eu^{\iu\varphi} \in \mathbb D$,
\[
u(z) = \frac1{2\pi}\int_0^{2\pi}
P_r(\varphi - t)\,g(\eu^{\iu t})\,\dd t .
\]
Alors $u$ est harmonique sur $\mathbb D$ et s'étend
continûment à $\bar{\mathbb D}$ avec valeurs au bord $g$~:
l'unique telle \omterm{prop:b3:conformal:harmonicholo}{fonction
harmonique}.\index{problème de Dirichlet}
\end{theorem}

\begin{proof}
\emph{Identités du noyau}~: en sommant deux séries
géométriques,
\[
\sum_{n\in\Z}r^{\abs n}\eu^{\iu n\theta}
= \operatorname{Re}\frac{1 + r\eu^{\iu\theta}}{1 -
r\eu^{\iu\theta}}
= \frac{1 - r^2}{\abs{1 - r\eu^{\iu\theta}}^2},
\]
qui est le quotient affiché~; la positivité est claire, et
$\frac1{2\pi}\int_0^{2\pi}P_r = 1$ (seul $n = 0$ survit).

\emph{Harmonicité}~: avec $z = r\eu^{\iu\varphi}$,
\[
u(z) = \operatorname{Re}\biggl[\frac1{2\pi}\int_0^{2\pi}
\frac{\eu^{\iu t} + z}{\eu^{\iu t} - z}\,
g(\eu^{\iu t})\,\dd t\biggr],
\]
(le noyau entre crochets a pour partie réelle $P_r(\varphi -
t)$~: calculer), et le crochet est \omterm{def:b3:holomorphic:holo}{holomorphe} en $z$ sur
$\mathbb D$ (l'\cref{exo:b3:holomorphic:7})~: $u$ est la
partie réelle d'une \omterm{def:b3:holomorphic:holo}{fonction holomorphe}, d'où harmonique.

\emph{Valeurs au bord}~: $P_r(\cdot)$ est une identité
approximative quand $r \to 1^-$~: masse $1$, et pour $\delta
\leq \abs\theta \leq \pi$, $P_r(\theta) \leq \frac{1 -
r^2}{1 - 2r\cos\delta + r^2} \to 0$ uniformément. Le
découpage standard (continuité de $g$ près de
$\eu^{\iu\varphi_0}$, bornitude ailleurs) donne
$u(r\eu^{\iu\varphi}) \to g(\eu^{\iu\varphi_0})$ quand
$r\eu^{\iu\varphi} \to \eu^{\iu\varphi_0}$, uniformément en
le point du bord~: l'extension est \omterm{def:b3:topology:continuity}{continue}. \emph{Unicité}~:
la différence de deux solutions est harmonique sur $\mathbb
D$, \omterm{def:b3:topology:continuity}{continue} sur l'\omterm{def:b3:topology:interior}{adhérence}, nulle sur le bord~: par le
principe du maximum (appliqué à $\pm$ la différence), elle
s'annule.
\end{proof}

\begin{omfigure}
\begin{tikzpicture}[scale=1.05]
  \begin{scope}
    \clip (-2.1,-2.1) rectangle (2.1,2.1);
    \foreach \c in {0.2, 0.5, 0.8, 1.1, 1.4}{
      \draw[omDef!60] plot[domain=-2.05:2.05, samples=80]
        (\x, {\c/(2*\x)});
      \draw[omDef!60] plot[domain=-2.05:2.05, samples=80]
        (\x, {-\c/(2*\x)});
      \draw[omThm!60, samples=80, domain=\c:2.9]
        plot ({sqrt((\x*\x-\c*\c)/1)/1}, {0}) ;
    }
    % hyperbolas x^2 - y^2 = c
    \foreach \c in {0.4, 1, 1.6}{
      \draw[omThm!70] plot[domain=-1.32:1.32, samples=60]
        ({cosh(\x)*sqrt(\c)}, {sinh(\x)*sqrt(\c)});
      \draw[omThm!70] plot[domain=-1.32:1.32, samples=60]
        ({-cosh(\x)*sqrt(\c)}, {sinh(\x)*sqrt(\c)});
    }
  \end{scope}
  \draw[->] (-2.2,0) -- (2.3,0);
  \draw[->] (0,-2.2) -- (0,2.3);
  \node[font=\small, omThm] at (1.75,-1.9)
    {$\operatorname{Re}z^2 = c$};
  \node[font=\small, omDef] at (-1.35,1.95)
    {$\operatorname{Im}z^2 = c$};
\end{tikzpicture}

{\small Les courbes de niveau de $\operatorname{Re}z^2 = x^2
- y^2$ (hyperboles rouges) et $\operatorname{Im}z^2 = 2xy$
(hyperboles bleues) se coupent à angle droit hors de $0$~:
une \omterm{def:b3:conformal:conformal}{application conforme} ($z \mapsto z^2$, où $z \neq 0$)
préserve l'orthogonalité des lignes de coordonnées. En $z =
0$, où la dérivée s'annule, les angles sont doublés à la
place.}
\end{omfigure}

\section{Exercices}

\begin{exercise}[$\star$]\label{exo:b3:conformal:1}
(a) Vérifier que l'\omterm{ex:b3:conformal:mobius}{application de Cayley} $\varphi(z) =
\frac{z - \iu}{z + \iu}$ est une bijection $\mathbb H \to
\mathbb D$ avec l'inverse indiqué, et calculer les images de
$\iu$, $0$, $1$, $\infty$ (limite).
(b) Montrer que $z \mapsto 1/z$ envoie cercles et droites sur
cercles et droites. \emph{(Écrire leur équation commune
$\alpha\abs z^2 + \bar\beta z + \beta\bar z + \gamma = 0$,
$\alpha, \gamma \in \R$.)} En déduire le même pour toutes
les applications de Möbius.
\end{exercise}

\begin{exercise}[$\star$]\label{exo:b3:conformal:2}
Soit $f \colon \mathbb D \to \mathbb D$ \omterm{def:b3:holomorphic:holo}{holomorphe}.
(a) Si $f(0) = 0$ et $f(a) = a$ pour un $a \neq 0$, montrer
$f = \mathrm{id}$.
(b) Si $f$ est un automorphisme avec deux points fixes
distincts dans $\mathbb D$, montrer $f = \mathrm{id}$
\emph{(conjuguer par un \omterm{thm:b3:conformal:autdisc}{facteur de Blaschke} pour se ramener
à (a))}.
\end{exercise}

\begin{exercise}[$\star\star$]\label{exo:b3:conformal:3}
(Schwarz--Pick) Pour $f\colon \mathbb D \to \mathbb D$
\omterm{def:b3:holomorphic:holo}{holomorphe}, prouver
\[
\frac{\abs{f'(z)}}{1 - \abs{f(z)}^2} \;\leq\;
\frac{1}{1 - \abs z^2}
\qquad (z \in \mathbb D),
\]
avec égalité (en un point, d'où partout) ssi $f \in
\operatorname{Aut}(\mathbb D)$.
\emph{(Appliquer Schwarz à $\varphi_{f(z)}\circ
f\circ\varphi_{-z}$.)} Interprétation~: les auto-applications
\omterm{def:b3:holomorphic:holo}{holomorphes} contractent la métrique hyperbolique.
\end{exercise}

\begin{exercise}[$\star\star$]\label{exo:b3:conformal:4}
Trouver des équivalences conformes explicites~:
(a) la bande $\{0 < \operatorname{Im}z < \pi\} \to \mathbb
H$~;
(b) le quadrant $\{\operatorname{Re}z > 0,
\operatorname{Im}z > 0\} \to \mathbb H$~;
(c) le demi-disque $\mathbb D\cap\mathbb H \to$ un quadrant,
puis $\to \mathbb H$~;
(d) $\mathbb D \to \mathbb D$ envoyant $\frac12$ en $0$ avec
dérivée positive là.
\end{exercise}

\begin{exercise}[$\star\star$]\label{exo:b3:conformal:5}
(a) Montrer qu'aucune \omterm{def:b3:conformal:conformal}{application conforme} $\C \to \mathbb D$
ou $\C \to \mathbb H$ n'existe, ni aucune $\mathbb D \to
\C$.
(b) Lesquelles des suivantes sont conformément équivalentes
à $\mathbb D$~? Justifier via le~\cref{thm:b3:conformal:rmt}
ou une obstruction~: un carré~;
$\C\setminus\intoc{-\infty}0$~; $\mathbb D\setminus\{0\}$~;
$\{1 < \abs z < 2\}$.
\emph{(Pour les deux dernières~: une image conforme du
disque percé s'étendrait sur la perforation par
le~\cref{thm:b3:residues:riemanncw}(1) --- développer
ceci.)}
\end{exercise}

\begin{exercise}[$\star\star$]\label{exo:b3:conformal:6}
Soit $\mathcal F = \{f \in \mathcal H(\mathbb D) : f(0) =
1,\ \operatorname{Re}f > 0\}$.
(a) Montrer que $\mathcal F$ est localement bornée.
\emph{(Composer avec l'application de type Cayley $w \mapsto
\frac{w - 1}{w + 1}$ envoyant le demi-plan droit sur
$\mathbb D$, et appliquer Schwarz.)}
(b) En déduire la \emph{borne de Herglotz}~: $\abs{f(z)}
\leq \frac{1 + \abs z}{1 - \abs z}$ pour $f \in \mathcal F$,
avec possibilités d'égalité.
\end{exercise}

\begin{exercise}[$\star\star\star$]\label{exo:b3:conformal:7}
Où la preuve du~\cref{thm:b3:conformal:rmt} utilise-t-elle
chaque hypothèse~? Tracer~: (i) \omterm{def:b3:groups:simple}{simple} connexité (deux
fois)~; (ii) $\Omega \neq \C$~; (iii) connexité. Puis
montrer que le théorème \emph{échoue} pour $\Omega = \C$ et
pour l'anneau, en localisant quelle étape de la preuve casse
dans chaque cas.
\end{exercise}

\begin{exercise}[$\star\star$]\label{exo:b3:conformal:8}
Résoudre le problème de Dirichlet sur $\mathbb D$ pour les
données au bord~: (a) $g(\eu^{\iu t}) = \cos t$~; (b)
$g(\eu^{\iu t}) = \cos^2 t$~; (c) $g = \mathbf
1_{\text{demi-cercle supérieur}}$ --- pour (c) calculer
$u(0)$ et interpréter via la propriété de la moyenne.
\emph{(Développer $g$ en série de Fourier et utiliser la
série de $P_r$~: $u(r\eu^{\iu\varphi}) = \sum_n c_n(g)
r^{\abs n}\eu^{\iu n\varphi}$.)}
\end{exercise}

\begin{exercise}[$\star\star\star$]\label{exo:b3:conformal:9}
(Harnack) Soit $u \geq 0$ harmonique sur $\mathbb D$.
Prouver, pour $\abs z = r < 1$~:
\[
\frac{1 - r}{1 + r}\,u(0) \;\leq\; u(z) \;\leq\;
\frac{1 + r}{1 - r}\,u(0)
\]
\emph{(borner le \omterm{thm:b3:conformal:poisson}{noyau de Poisson} entre $\frac{1-r}{1+r}$ et
$\frac{1+r}{1-r}$~; appliquer la \omterm{def:b3:representations:rep}{représentation} sur des
disques légèrement plus petits et passer à la limite)}. En
déduire~: une \omterm{prop:b3:conformal:harmonicholo}{fonction harmonique} sur $\C$ bornée
inférieurement est constante.
\end{exercise}

\begin{exercise}[$\star\star$]\label{exo:b3:conformal:10}
En utilisant l'invariance conforme de l'harmonicité ($u\circ
f$ est harmonique quand $u$ est harmonique et $f$ \omterm{def:b3:holomorphic:holo}{holomorphe}
--- le prouver via la~\cref{prop:b3:conformal:harmonicholo}
localement), résoudre le problème de Dirichlet sur le
demi-plan supérieur avec données au bord $\mathbf
1_{\intoo{-\infty}0}$~: montrer que
\[
u(x + \iu y) = \frac1\pi\,\arg(x + \iu y)
\qquad (\arg \in \intoo0\pi \text{ sur } \mathbb H)
\]
est harmonique sur $\mathbb H$ (partie imaginaire d'un
\omterm{def:b3:conformal:simplyconnected}{logarithme holomorphe}) avec les limites au bord requises en
tout $x \neq 0$, et le transporter au disque par Cayley pour
retrouver l'\cref{exo:b3:conformal:8}(c).
\end{exercise}

\begin{exercise}[$\star\star$]\label{exo:b3:conformal:11}
(Points fixes et itération dans le disque) Soit
$f\colon\mathbb D \to \mathbb D$ \omterm{def:b3:holomorphic:holo}{holomorphe}.
(a) Montrer que si $f$ a \emph{deux} points fixes distincts,
alors $f = \mathrm{id}$ \emph{(amener l'un en $0$ par un
automorphisme et appliquer le cas d'égalité de Schwarz)}.
(b) Supposer $f(0) = 0$ et $f$ n'est pas une rotation.
Montrer que les itérées $f^{\circ n} \to 0$ uniformément sur
tout compact $\bar D(0, r)$, $r < 1$ \emph{(Schwarz donne
$\abs{f(z)} \leq c_r\abs z$ sur $\bar D(0,r)$ avec $c_r < 1$
--- justifier cette constante stricte via le principe du
maximum appliqué à $f(z)/z$)}.
(c) Illustrer avec $f(z) = \frac{z^2 + z}2$~: points fixes,
et le taux de convergence de l'\omterm{def:b3:groups:action}{orbite} de $z_0 =
\frac12$.
\end{exercise}

\begin{exercise}[$\star\star$]\label{exo:b3:conformal:12}
(Conjuguées harmoniques, concrètement) Soit $u(x, y) = x^3
- 3xy^2 + 2y$.
(a) Vérifier que $u$ est harmonique sur $\R^2$, et trouver
toutes les conjuguées harmoniques $v$ (c'est-à-dire $u + \iu
v$ \omterm{def:b3:holomorphic:holo}{holomorphe}) en intégrant les \omterm{prop:b3:holomorphic:cauchyriemann}{équations de
Cauchy--Riemann}~; identifier $f(z) = u + \iu v$ comme un
polynôme en $z$.
(b) Montrer que sur un \emph{étoilé} \omterm{def:b3:topology:topology}{ouvert}, toute \omterm{prop:b3:conformal:harmonicholo}{fonction
harmonique} admet une conjuguée harmonique, unique à une
constante additive près \emph{(la $1$-forme $-u_y\,\dd x +
u_x\,\dd y$ est fermée~; la machinerie de primitive
du~\cref{thm:b3:holomorphic:cauchy}, ou le lemme de Poincaré
du~\cref{ch:b3:forms})}.
(c) Donner le contre-exemple standard sur $\C^*$~: $u =
\ln\abs z$ n'a pas de conjuguée globale --- relier à la
forme angulaire et au \omterm{def:b3:holomorphic:index}{nombre de tours}
(le~\cref{ch:b3:forms}).
\end{exercise}

\section{Problème~: le théorème de l'aire et le théorème du
quart de Koebe}

\begin{problem}[{Problème de week-end --- combien une
application univalente doit-elle couvrir~?}]\label{pb:b3:conformal:1}
Une fonction \emph{univalente} est une injection \omterm{def:b3:holomorphic:holo}{holomorphe}.
Les fonctions univalentes normalisées sur le disque,
\[
\mathcal S = \bigl\{f \in \mathcal H(\mathbb D) \text{
injective},\ f(z) = z + a_2z^2 + a_3z^3 + \cdots\bigr\},
\]
sont rigidement contraintes~: nous prouvons l'inégalité de
Bieberbach $\abs{a_2} \leq 2$ et en déduisons le
\emph{théorème du quart de Koebe}~: l'image de toute $f \in
\mathcal S$ contient le disque $D(0, \frac14)$ --- la
constante universelle nette de la géométrie
conforme.\index{théorème du quart de Koebe}\index{théorème de
l'aire}

\medskip
\noindent\textbf{Partie~I --- Le théorème de l'aire.} Soit
$g(w) = w + b_0 + \frac{b_1}w + \frac{b_2}{w^2} + \cdots$
\omterm{def:b3:holomorphic:holo}{holomorphe} et \emph{injective} sur $\{\abs w > 1\}$.
\begin{enumerate}
  \item Pour $\rho > 1$, soit $A_\rho$ l'aire (\omterm{def:b3:measure:lebesgueouter}{mesure de
        Lebesgue}) de l'ensemble compact $K_\rho =
        \C\setminus g(\{\abs w > \rho\})$, la région
        encerclée par la courbe de Jordan lisse $g(C_\rho)$.
        En utilisant la formule d'aire du théorème de
        Green--Riemann du volume de deuxième année ---
        l'aire encerclée est $\frac1{2\iu}
        \oint\bar\zeta\,\dd\zeta$ le long du bord orienté
        positivement --- montrer que
        \[
        A_\rho = \frac{1}{2\iu}\int_{C_\rho}
        \overline{g(w)}\,g'(w)\,\dd w
        = \pi\Bigl(\rho^2 -
        \sum_{n\geq1}n\,\abs{b_n}^2\rho^{-2n}\Bigr) :
        \]
        substituer la \omterm{thm:b3:residues:laurent}{série de Laurent} de $\bar g$ et $g'$
        sur $C_\rho$ et intégrer terme à terme (convergence
        normale~; seuls les produits de fréquence nulle
        survivent).
  \item Faire $\rho \to 1^+$ et conclure le \emph{théorème de
        l'aire}~:
        \[
        \sum_{n\geq1}n\,\abs{b_n}^2 \;\leq\; 1 .
        \]
        En particulier $\abs{b_1} \leq 1$. Quand a-t-on
        $\abs{b_1} = 1$~?
\end{enumerate}

\medskip
\noindent\textbf{Partie~II --- Bieberbach $\abs{a_2} \leq
2$.} Soit $f = z + a_2z^2 + \cdots \in \mathcal S$.
\begin{enumerate}[resume]
  \item Montrer que $f(z^2)/z^2$ est \omterm{def:b3:holomorphic:holo}{holomorphe} et sans zéro
        sur $\mathbb D$, et admet une racine carrée
        \omterm{def:b3:holomorphic:holo}{holomorphe} $\varphi$ avec $\varphi(0) = 1$~; poser
        $h(z) = z\varphi(z^2)$, de sorte que $h(z)^2 =
        f(z^2)$. Montrer que $h$ est une fonction univalente
        \emph{impaire} sur $\mathbb D$ de développement $h(z)
        = z + \frac{a_2}2z^3 + \cdots$. \emph{(Injectivité~:
        $h(z)^2 = h(z')^2$ force $z^2 = z'^2$~; utiliser
        l'imparité pour finir.)}
  \item Appliquer le théorème de l'aire à $g(w) = 1/h(1/w) =
        w - \frac{a_2}{2w} + \cdots$ sur $\{\abs w > 1\}$
        (vérifier l'univalence et le développement), et
        conclure $\abs{a_2} \leq 2$.
  \item Montrer que la \emph{fonction de Koebe}
        \[
        k(z) = \frac{z}{(1 - z)^2} = \sum_{n\geq1}n\,z^n
        \]
        appartient à $\mathcal S$, a $a_2 = 2$, et envoie
        $\mathbb D$ sur $\C\setminus\intoc{-\infty}{-\frac14}$
        \emph{(écrire $k = \frac14\bigl[\bigl(\frac{1+z}{1 -
        z}\bigr)^2 - 1\bigr]$ et suivre les images)}~:
        toutes les inégalités à venir sont nettes.
\end{enumerate}

\medskip
\noindent\textbf{Partie~III --- Le théorème du quart.}
\begin{enumerate}[resume]
  \item Soit $f \in \mathcal S$ et $c \notin f(\mathbb D)$.
        Montrer que
        \[
        F(z) = \frac{c\,f(z)}{c - f(z)}
        \]
        appartient à $\mathcal S$, et calculer son second
        coefficient~: $A_2 = a_2 + \frac1c$.
  \item Appliquer Bieberbach à $f$ et à $F$~: conclure
        $\abs{\frac1c} \leq \abs{A_2} + \abs{a_2} \leq 4$,
        c'est-à-dire $\abs c \geq \frac14$. \emph{Toute
        valeur omise a un \omterm{def:b3:modules:module}{module} $\geq \frac14$}~:
        $f(\mathbb D) \supseteq D(0, \frac14)$ --- théorème
        du quart de Koebe. Vérifier la netteté sur la
        fonction de Koebe.
  \item En déduire une estimation quantitative d'application
        de Riemann~: si $\varphi \colon \Omega \to \mathbb
        D$ est l'application de Riemann
        du~\cref{thm:b3:conformal:rmt} en $z_0$, alors
        \[
        \frac{d\bigl(z_0, \partial\Omega\bigr)}{4}
        \;\leq\; \frac1{\varphi'(z_0)} \;\leq\;
        4\,d\bigl(z_0, \partial\Omega\bigr)
        \]
        --- prouver au moins l'inégalité de gauche en
        appliquant Koebe à $\varphi^{-1}$ convenablement
        normalisée, et celle de droite par Schwarz appliqué
        à $\varphi$ sur le disque $D(z_0, d) \subseteq
        \Omega$.
\end{enumerate}

\medskip
\noindent\textbf{Partie~IV --- Perspective.}
\begin{enumerate}[resume]
  \item Bieberbach a conjecturé (1916) $\abs{a_n} \leq n$
        pour tout $n$, avec égalité seulement pour les
        rotations de la fonction de Koebe~; de Branges l'a
        prouvé en 1985. Vérifier la conjecture à la main
        pour la fonction de Koebe et ses rotations
        $\eu^{-\iu\theta}k(\eu^{\iu\theta}z)$. Puis pousser
        le développement de la question~4 d'un terme~: en
        écrivant $h(z) = z + \frac{a_2}2z^3 + c_5z^5 +
        \cdots$, montrer $c_5 = \frac{a_3}2 -
        \frac{a_2^2}8$ et
        \[
        g(w) = w - \frac{a_2}{2}\,w^{-1} +
        \Bigl(\frac{3a_2^2}8 - \frac{a_3}2\Bigr)w^{-3} +
        \cdots ,
        \]
        donc le théorème de l'aire donne l'inégalité affinée
        $\bigl|\frac{a_2}2\bigr|^2 + 3\bigl|\frac{3a_2^2}8 -
        \frac{a_3}2\bigr|^2 \leq 1$. Vérifier sur la fonction
        de Koebe ($a_2 = 2$, $a_3 = 3$).
\end{enumerate}

\medskip
\noindent\textbf{Partie~V --- Le théorème de distorsion.}
L'inégalité de Bieberbach, transportée autour du disque par
les automorphismes, contrôle $f'$ partout. Fixer $f \in
\mathcal S$.
\begin{enumerate}[resume]
  \item (Transformée de Koebe) Pour $z_0 \in \mathbb D$
        soit $\varphi(z) = \frac{z + z_0}{1 + \bar
        z_0z}$, un automorphisme du disque
        (le~\cref{thm:b3:conformal:autdisc}) avec
        $\varphi(0) = z_0$. Montrer que
        \[
        F(z) = \frac{f(\varphi(z)) - f(z_0)}
        {f'(z_0)\,\bigl(1 - \abs{z_0}^2\bigr)}
        \]
        appartient à $\mathcal S$ \emph{(l'univalence est
        héritée~; calculer $\varphi'(0) = 1 - \abs{z_0}^2$
        et vérifier la normalisation~; rappeler $f' \neq 0$
        pour $f$ injective,
        la~\cref{def:b3:conformal:conformal})}.
  \item Calculer le second coefficient $A_2 =
        \frac12F''(0)$ de $F$~:
        \[
        A_2 = \frac12\Bigl[\bigl(1 - \abs{z_0}^2\bigr)
        \frac{f''(z_0)}{f'(z_0)} - 2\bar z_0\Bigr] ,
        \]
        et déduire de Bieberbach (question~4), pour $z =
        r\eu^{\iu\theta}$, l'\emph{inégalité
        fondamentale}~:
        \[
        \Bigl|\,z\,\frac{f''(z)}{f'(z)} -
        \frac{2r^2}{1 - r^2}\Bigr|
        \;\leq\; \frac{4r}{1 - r^2} .
        \]
  \item Extraire la partie réelle~:
        \[
        \frac{2r^2 - 4r}{1 - r^2} \;\leq\;
        \operatorname{Re}\Bigl(z\,\frac{f''(z)}{f'(z)}\Bigr)
        \;\leq\; \frac{2r^2 + 4r}{1 - r^2} .
        \]
  \item Montrer que $\frac{\dd}{\dd t}\log\bigl|
        f'(t\eu^{\iu\theta})\bigr| = \frac1t
        \operatorname{Re}\bigl(z\frac{f''(z)}{f'(z)}\bigr)$
        en $z = t\eu^{\iu\theta}$ \emph{(pour une fonction
        $\mathcal C^1$ non nulle $g$ d'une variable réelle,
        $\frac{\dd}{\dd t}\log\abs g =
        \operatorname{Re}(g'/g)$)}, et intégrer les bornes
        de la question~12 le long du rayon pour obtenir le
        \emph{théorème de distorsion}~:
        \[
        \frac{1 - r}{(1 + r)^3} \;\leq\; \abs{f'(z)}
        \;\leq\; \frac{1 + r}{(1 - r)^3},
        \qquad \abs z = r .
        \]
  \item En déduire le \emph{théorème de croissance}~:
        \[
        \frac{r}{(1 + r)^2} \;\leq\; \abs{f(z)} \;\leq\;
        \frac{r}{(1 - r)^2},
        \qquad \abs z = r
        \]
        \emph{(borne supérieure~: intégrer $f'$ sur le
        segment $[0, z]$~; borne inférieure~: si $\abs{f(z)}
        < \frac14$, le segment $[0, f(z)]$ est dans
        $f(\mathbb D)$ par la question~7~; le ramener par
        $f^{-1}$ --- \omterm{def:b3:holomorphic:holo}{holomorphe} par
        le~\cref{cor:b3:residues:openmapping} --- et borner
        $\abs{f(z)} = \int_\gamma\abs{f'(\zeta)}
        \,\abs{\dd\zeta} \geq \int_0^r\frac{1 - t}{(1 +
        t)^3}\,\dd t$, en utilisant $\abs{\dd\zeta} \geq
        \dd\abs\zeta$)}.
  \item Vérifier que la fonction de Koebe réalise l'égalité
        dans les quatre bornes, en $z = r$ pour les
        supérieures et $z = -r$ pour les inférieures~: $k$
        est simultanément le membre le plus expansif et, à
        l'antipode, le plus contractant de $\mathcal S$.
\end{enumerate}

\medskip
\noindent\textbf{Partie~VI --- Rigidité extrémale.} Dans
chaque inégalité jusqu'ici, l'égalité identifie la fonction
de Koebe à rotation près. Nous le prouvons, puis récoltons.
\begin{enumerate}[resume]
  \item Supposer $f \in \mathcal S$ avec $\abs{a_2} = 2$.
        Suivre l'égalité à travers les questions~2--4~: le
        théorème de l'aire force $g(w) = w + b_0 +
        \eu^{\iu\alpha}/w$~; l'imparité de $h$ rend $g$
        impaire, donc $b_0 = 0$~; inverser pour trouver $h$,
        puis $f$, et conclure que
        \[
        f(z) = \eu^{-\iu\theta}k\bigl(\eu^{\iu\theta}z\bigr)
        \quad\text{avec } \eu^{\iu\theta} =
        -\eu^{\iu\alpha} :
        \]
        les rotations de la fonction de Koebe sont les seuls
        membres de $\mathcal S$ avec $\abs{a_2} = 2$.
  \item Montrer que si $f \in \mathcal S$ omet une valeur $c$
        avec $\abs c = \frac14$ exactement, alors $f$ est
        une rotation de la fonction de Koebe, et identifier
        $c = -\eu^{-\iu\theta}/4$ \emph{(suivre l'égalité à
        travers la chaîne de la question~7 $4 = \abs{1/c} =
        \abs{A_2 - a_2} \leq \abs{A_2} + \abs{a_2} \leq
        4$)}~: la constante du théorème du quart n'est
        atteinte que par la famille extrémale.
  \item (Coefficients à bon marché) Combiner le théorème de
        croissance avec les \omterm{thm:b3:holomorphic:analytic}{estimations de Cauchy}
        (le~\cref{thm:b3:holomorphic:analytic}) sur le
        cercle $\abs z = 1 - \frac1n$ pour prouver
        \[
        \abs{a_n} \;\leq\; \eu\,n^2 \qquad (n \geq 2) .
        \]
        (De Branges, 1985~: $\abs{a_n} \leq n$~; le facteur
        $\eu n$ est le prix d'outils élémentaires.)
  \item (Recouvrement de sous-disques) Montrer que pour tout
        $0 < r < 1$,
        \[
        f\bigl(D(0, r)\bigr) \supseteq
        D\Bigl(0, \frac{r}{(1 + r)^2}\Bigr),
        \]
        net pour la fonction de Koebe, et retrouver le
        théorème du quart comme $r \to 1^-$. \emph{(Les
        points du bord de l'image ouverte $f(D(0,r))$ sont
        sur $f(\partial D(0,r))$, d'où de \omterm{def:b3:modules:module}{module} $\geq
        \frac{r}{(1+r)^2}$ par la question~14~; un segment
        de $0$ à un point manqué de plus petit \omterm{def:b3:modules:module}{module}
        devrait croiser ce bord.)}
  \item (Koebe en tout point) Soit $f$ univalente sur
        $\mathbb D$, pas nécessairement normalisée, et $z_0
        \in \mathbb D$. Prouver
        \[
        \tfrac14\bigl(1 - \abs{z_0}^2\bigr)\abs{f'(z_0)}
        \;\leq\;
        d\bigl(f(z_0), \partial f(\mathbb D)\bigr)
        \;\leq\;
        \bigl(1 - \abs{z_0}^2\bigr)\abs{f'(z_0)}
        \]
        \emph{(gauche~: théorème du quart appliqué à la
        transformée de Koebe de la question~10~; droite~:
        Schwarz (le~\cref{thm:b3:conformal:schwarz})
        appliqué à $\psi^{-1}\circ\hat g$, où $\hat g(w) =
        f^{-1}\bigl(f(z_0) + dw\bigr)$, $d$ la distance, et
        $\psi$ un automorphisme du disque envoyant $0$ en
        $z_0$)}. Pourquoi $\partial f(\mathbb D)$ est-il non
        vide~?
  \item Vérifier la question~20 sur $f = k$ en $z_0 = r \in
        \intoo01$~: calculer $d\bigl(k(r), \partial
        k(\mathbb D)\bigr) = \frac{(1+r)^2}{4(1-r)^2}$ et
        vérifier que l'inégalité de gauche est une
        \emph{égalité}~: la fonction de Koebe sature son
        propre théorème en tout point de $\intoo01$.
  \item (La morale) En dix lignes~: quel principe unique
        sous-tend le théorème de l'aire, et comment
        Bieberbach, le théorème du quart, la distorsion, la
        croissance et le recouvrement en découlent tous~?
        Comparer avec le monde de Schwarz--Pick
        du~\cref{thm:b3:conformal:schwarz}~: dans les deux,
        une inégalité intérieure rigidifie toute la
        géométrie, et les extrémales sont uniques à
        rotation près.

\end{enumerate}
\medskip
\noindent\textbf{Partie~VII --- Compacité, inverses, et un
contrôle de réalité.}
\begin{enumerate}[resume]
  \item Montrer que la classe $\mathcal S$ est
        \emph{compacte} pour la convergence localement
        uniforme~: elle est localement bornée par le
        théorème de croissance, d'où normale
        (le~\cref{thm:b3:conformal:montel})~; et une limite
        localement uniforme de membres de $\mathcal S$ est
        encore dans $\mathcal S$ \emph{(les normalisations
        passent à la limite par convergence de Weierstrass
        des dérivées~; l'injectivité survit par Hurwitz,
        l'\cref{exo:b3:residues:8}, la limite étant non
        constante)}. Pourquoi cela importe-t-il pour les
        problèmes extrémaux comme celui de Bieberbach~?
  \item Pour $f \in \mathcal S$, soit $g = f^{-1}$, définie
        près de $0$. Montrer $g(w) = w - a_2w^2 + O(w^3)$,
        donc le second coefficient de l'inverse obéit à la
        même borne nette $\abs{A_2} = \abs{a_2} \leq 2$,
        avec égalité exactement pour les fonctions de Koebe
        tournées.
  \item Déterminer pour quels $a \in \C$ le polynôme $f(z)
        = z + az^2$ appartient à $\mathcal S$~: montrer que
        $f$ est injective sur $\mathbb D$ ssi $\abs a \leq
        \frac12$ \emph{(factoriser $f(z_1) - f(z_2)$)}.
        Conclure~: pour les polynômes de \omterm{def:b3:galois:extension}{degré} deux la vraie
        borne de coefficient est $\frac12$, quatre fois plus
        petite que le $2$ de Bieberbach --- les extrémales
        de $\mathcal S$ sont authentiquement des objets
        transcendants, et aucun polynôme ne s'en approche.
\end{enumerate}
\end{problem}
